\section{Variational Streaming Flow: Method and Derivation}
\label{sec:method}





%\paragraph{Continuous-time Variational Streaming Flow (VSF).} 


Motivated to address this limitation, we augment the physical-time velocity field
with a latent variable. 
Let $ \mathbf{Z}
    \sim
    p(\vz)
    :=
    \mathcal N
    \bigl(
        \vz;
        \mathbf{0},
        \mathbf{I}_{d_z}
    \bigr)$  be a latent variable.  Conditional on $(\vh,\mathbf{Z})$, VSF evolves according to
\begin{tcolorbox}[
  colback=parchment,
  colframe=sandborder,
  boxrule=0.3pt,
  left=8pt,
  right=8pt,
  top=3pt,
  bottom=3pt
]
\begin{equation}
\frac{d\widehat{\vx}_s}{ds}
=
\vv_{\vtheta}\!\left(
    \widehat{\vx}_s,s,\vh,\vz
\right), \qquad \widehat{\vx}_0=\vx_m, 
\qquad s\in[0,1].
\label{eq:vsf_dynamics}
\end{equation}
\end{tcolorbox}


During training, a realized future trajectory $\boldsymbol{\xi}$ is available.
For a sampled physical time $s$, define the local state and velocity as
$\vx_s:=\boldsymbol{\xi}(s),$ $\vu_s:=\dot{\boldsymbol{\xi}}(s).$
We infer the latent locally using:
\begin{equation}
q_{\boldsymbol{\phi}}
\bigl(
\vz
\mid
\vu_s,
\vx_s,
s,
\vh
\bigr)
=
\mathcal N
\left(
\vz;
\boldsymbol{\mu}_{\boldsymbol{\phi}}
(\vu_s,\vx_s,s,\vh),
\operatorname{diag}
\left[
\boldsymbol{\sigma}_{\boldsymbol{\phi}}^2
(\vu_s,\vx_s,s,\vh)
\right]
\right).
\label{eq:vsf_variational_posterior}
\end{equation}
The variational posterior $q_{\boldsymbol{\phi}}$, parameterized by the
recognition network, provides a  
local latent posterior induced by the conditional velocity model.
At inference, the future state and
its target velocity are unavailable, so the recognition network is discarded.
Instead, $\mathbf{Z}\sim p(\vz)$ is sampled once and held fixed throughout
the rollout. In \cref{sec:ablation-studies}, we empirically show that fixing
$\mathbf{Z}$ throughout the rollout improves long-horizon predictive accuracy
relative to resampling it at every step.







\paragraph{Discrete observations as a realization of the continuous model.} 
The preceding formulation treats physical time continuously.  In applications,
however, the future is observed only at discrete times.  For simplicity, we
assume equally spaced observation times (for irregular sampling, the
same construction applies with normalized timestamps and interval-dependent
scaling). 
 
Consider a training window $\vx_{a:T}$ with forecast origin
$m$, and let $P := T-m$ denote the number of future transitions. For $k = 0, \cdots, P$, define the grid $s_k := k/P$, and, for $k=0, \cdots, P-1$, $ I_k:=[s_k,s_{k+1}), \alpha_s:=P(s-s_k)$ for $s\in I_k$. Since the continuous path between two observations is unavailable, we
construct the piecewise-linear surrogate
$\boldsymbol{\xi}^{\Delta}:[0,1]\to\mathcal X$,
\begin{equation}
    \boldsymbol{\xi}^{\Delta}(s)
    :=
    (1-\alpha_s)\vx_{m+k}
    +
    \alpha_s\vx_{m+k+1},
    \qquad
    s\in I_k,
    \label{eq:physical_time_interpolation}
\end{equation}
with
$\boldsymbol{\xi}^{\Delta}(0):=\vx_m$
and
$\boldsymbol{\xi}^{\Delta}(1):=\vx_T$. Its derivative is a local velocity which  is constant on each interval: 
\begin{equation}
    \dot{\boldsymbol{\xi}}^{\Delta}(s)
    =
    P
    \bigl(
        \vx_{m+k+1}-\vx_{m+k}
    \bigr)
    =:
    \vu_k,
    \qquad
    s\in I_k.
    \label{eq:physical_time_velocity}
\end{equation}


We next relate the continuous and discrete formulations. Let $K \sim\operatorname{Unif}\{0,\ldots,P-1\}$, $A\sim\operatorname{Unif}[0,1]$ be independent and set $S := (K+A)/P$. Then $S \sim \operatorname{Unif}[0,1]$ and 
\begin{equation}
    \boldsymbol{\xi}^{\Delta}(S)
    =
    (1-A)\vx_{m+K}
    +
    A\vx_{m+K+1},
 \qquad 
    \dot{\boldsymbol{\xi}}^{\Delta}(S)
    =
    P
    \bigl(
        \vx_{m+K+1}-\vx_{m+K}
    \bigr).
    \label{eq:sampled_velocity}
\end{equation}
Thus,  sampling a transition index $K$ and a within-transition interpolation coordinate $A$ is equivalent to sampling a uniformly distributed
physical time $S$ along the surrogate path. For concrete method instantiation, we adopt the discrete-time formulation from now on.



\paragraph{Variational approach of VSF.} Given the local state and velocity target obtained from the discrete surrogate,
we next formulate the training objective.
%The generative dynamics use a single latent variable that is held fixed throughout the rollout, while during training the latent is inferred locally from the sampled physical-time transition.  
We first consider a Gaussian likelihood for the local velocity, which yields a
variational MSE reconstruction. 
Although the latent variable allows this model to
represent multiple velocity modes, the reconstruction term acts on individual
latent--target pairs. We therefore replace it with a distributional objective
that compares velocities generated from multiple latent samples with the
target conditional velocity distribution.



%The generative dynamics use a single latent variable that is held fixed throughout the rollout.  During training, however, the latent is inferred locally from the sampled physical-time transition.


For a sampled physical time $S\sim\operatorname{Unif}[0,1]$, we evaluate
the piecewise-linear interpolation of the discrete observations to define
the local context and velocity target:
% The discrete surrogate provides the local state and velocity used by the
% variational encoder.  For a sampled physical time
% $S\sim\operatorname{Unif}[0,1]$, define
\begin{equation}
\vc_S
:=
\bigl(
\boldsymbol{\xi}^{\Delta}(S),
S,
\vh
\bigr),
\qquad
\vu_S
:=
\dot{\boldsymbol{\xi}}^{\Delta}(S).
\label{eq:clean_context_target}
\end{equation}
Thus the continuous-time recognition model specializes to
\begin{equation}
q_{\boldsymbol{\phi}}
\bigl(
\vz
\mid
\vu_S,
\vc_S
\bigr)
=
q_{\boldsymbol{\phi}}
\bigl(
\vz
\mid
\dot{\boldsymbol{\xi}}^{\Delta}(S),
\boldsymbol{\xi}^{\Delta}(S),
S,
\vh
\bigr).
\label{eq:discrete_local_posterior}
\end{equation}
The latent is therefore inferred locally during training, while at inference a
single draw $\mathbf{Z}\sim p(\vz)$ is retained throughout the
rollout.


A convenient variational starting point is the auxiliary Gaussian conditional likelihood
\begin{equation}
p_{\vtheta}
\bigl(
\vu_S
\mid
\vc_S,
\vz
\bigr)
=
\mathcal N
\left(
\vu_S;
\vv_{\vtheta}
\bigl(
\boldsymbol{\xi}^{\Delta}(S),
S,
\vh,
\vz
\bigr),
\sigma_u^2\mathbf{I}_d
\right),
\qquad
\sigma_u>0.
\label{eq:gaussian_velocity_likelihood}
\end{equation}
This likelihood models the local physical-time velocity rather than the full
trajectory law induced by the VSF ODE.  Marginalizing $\vz$ gives
\begin{equation}
p_{\vtheta}
(\vu_S\mid\vc_S)
=
\int
p_{\vtheta}
(\vu_S\mid\vc_S,\vz)
p(\vz)\,d\vz,
\label{eq:vsf_velocity_evidence}
\end{equation}
and the corresponding local ELBO is
\begin{align}
\log
p_{\vtheta}
(\vu_S\mid\vc_S)
\geq{}&
\mathbb E_{
\mathbf{Z}\sim
q_{\boldsymbol{\phi}}
(\cdot\mid\vu_S,\vc_S)
}
\left[
\log
p_{\vtheta}
(\vu_S\mid\vc_S,\mathbf{Z})
\right]
-
\operatorname{KL}
\left(
q_{\boldsymbol{\phi}}
(\vz\mid\vu_S,\vc_S)
\ \Vert \ 
p(\vz)
\right).
\label{eq:vsf_elbo}
\end{align}
For the Gaussian likelihood, the negative reconstruction term reduces, up to
an additive constant, to
\begin{align}
&-
\mathbb E_{
\mathbf{Z}\sim
q_{\boldsymbol{\phi}}
(\cdot\mid\vu_S,\vc_S)
}
\left[
\log
p_{\vtheta}
(\vu_S\mid\vc_S,\mathbf{Z})
\right]
\notag\\
&\qquad=
\frac{1}{2\sigma_u^2}
\mathbb E_{
\mathbf{Z}\sim
q_{\boldsymbol{\phi}}
(\cdot\mid\vu_S,\vc_S)
}
\left[
\left\lVert
\vv_{\vtheta}
\bigl(
\boldsymbol{\xi}^{\Delta}(S),
S,
\vh,
\mathbf{Z}
\bigr)
-
\vu_S
\right\rVert_2^2
\right]
+
\mathrm{const}.
\label{eq:vsf_variational_mse}
\end{align}
Thus the ELBO gives a latent-conditioned MSE reconstruction together
with prior regularization; see App.~\ref{sec:clean_velocity_elbo} for derivation of the corresponding  ELBO.

Let $(\mathbf{C},\mathbf{U})$ denote a generic random pair distributed as
$(\vc_S,\vu_S)$ under the training data and physical-time
sampling procedure, and let
\begin{equation}
\vz\mid(\mathbf{U},\mathbf{C})
\sim
q_{\boldsymbol{\phi}}(\cdot\mid\mathbf{U},\mathbf{C}).
\label{eq:generic_posterior_pair}
\end{equation}
For a fixed variational posterior $q_{\boldsymbol{\phi}}$, the population
minimizer of the MSE reconstruction  is
\begin{equation}
\vv^*(\vc,\vz)
=
\mathbb{E}
\left[
\mathbf{U}
\,\middle|\,
\mathbf{C}=\vc,
\mathbf{Z}=\vz
\right].
\label{eq:generic_mse_conditional_mean}
\end{equation}
Hence, when the latent representation does not fully distinguish different
future continuations, their residual local velocity ambiguity is averaged.
During training, such incomplete separation may arise from
overlapping branch-specific posteriors, imperfect  inference, or
posterior collapse.

To retain information about the distribution of plausible latent-conditioned
velocities, draw
\begin{equation}
\mathbf{Z}^{(i)}
\stackrel{\mathrm{iid}}{\sim}
q_{\boldsymbol{\phi}}
\bigl(
\vz
\mid
\vu_S,
\vc_S
\bigr),
\qquad
i=1,\ldots,M,
\label{eq:posterior_latent_samples}
\end{equation}
and define
\begin{equation}
\widehat{\vu}_S^{(i)}
:=
\vv_{\vtheta}
\bigl(
\boldsymbol{\xi}^{\Delta}(S),
S,
\vh,
\mathbf{Z}^{(i)}
\bigr).
\label{eq:posterior_velocity_samples}
\end{equation}
We use the generalized energy reconstruction
\begin{equation}
\mathcal L_{\mathrm{ES}}(S)
:=
\frac{1}{M}
\sum_{i=1}^M
\left\lVert
\widehat{\vu}_S^{(i)}
-
\vu_S
\right\rVert_2^\alpha
-
\frac{1}{2M^2}
\sum_{i,j=1}^M
\left\lVert
\widehat{\vu}_S^{(i)}
-
\widehat{\vu}_S^{(j)}
\right\rVert_2^\alpha,
\qquad
0<\alpha<2.
\label{eq:vsf_energy_score}
\end{equation}
For $\alpha=1$, this has the usual energy score form and reduces to the
CRPS in 1D~\citep{gneiting2007strictly}.
The first term penalizes discrepancy from the observed physical-time velocity,
while the second discourages unnecessary contraction of the  posterior
ensemble. Under aggregate posterior--prior alignment,
App.~\ref{sec:population_energy_control} shows that the expected
population energy reconstruction upper-bounds half the energy distance
between the prior-predictive and true conditional velocity laws.


The above discussions motivate the practical VSF objective:
\begin{tcolorbox}[
  colback=parchment,
  colframe=sandborder,
  boxrule=0.3pt,
  left=8pt,
  right=8pt,
  top=3pt,
  bottom=3pt
]
\begin{equation}
\mathcal L_{\mathrm{VSF}}
=
\mathbb E
\left[
\mathcal L_{\mathrm{ES}}(S)
+
\beta_{\mathrm{KL}}
\operatorname{KL}
\left(
q_{\boldsymbol{\phi}}
(\vz\mid\vu_S,\vc_S)
\,\Vert\,
p(\vz)
\right)
\right],
\label{eq:vsf_training_objective}
\end{equation}
\end{tcolorbox}
where the expectation is over training trajectories and sampled physical times. Note that once the Gaussian reconstruction is replaced by
\eqref{eq:vsf_energy_score}, this objective is no longer an ELBO; rather, it
is a KL-regularized distributional reconstruction motivated by the same local
latent model.
The complete batched training procedure and inference procedure are summarized in
\cref{alg:vsf-training}-\cref{alg:vsf-inference}.
%




\paragraph{Numerical rollout and the forecast.}
At inference, we draw one
$\mathbf{Z}\sim p(\vz)$
for each forecast member and retain it throughout the rollout.  Let
$\Psi_{\vtheta,k}^{\mathbf{Z},\vh}$ denote the numerical
flow map obtained by integrating \eqref{eq:vsf_dynamics} over
$I_k$.  The resulting discrete forecast satisfies
\begin{equation}
\widehat{\vx}_{m+k+1}
=
\Psi_{\vtheta,k}^{\mathbf{Z},\vh}
\bigl(
\widehat{\vx}_{m+k}
\bigr),
\qquad
k=0,\ldots,P-1,
\label{eq:vsf_discrete_flow_map}
\end{equation}
with $\widehat{\vx}_m=\vx_m$.  Thus the discrete forecast is
obtained by sampling the continuous physical-time trajectory on the
observation grid.
For a single Euler step over $I_k$,
\begin{equation}
\widehat{\vx}_{m+k+1}
=
\widehat{\vx}_{m+k}
+
\frac{1}{P}
\vv_{\vtheta}
\bigl(
\widehat{\vx}_{m+k},
s_k,
\vh,
\vz
\bigr).
\label{eq:vsf_euler_realization}
\end{equation}
%The factor $P$ in \eqref{eq:physical_time_velocity} is therefore precisely the normalization required when each observation interval has width $1/P$ in normalized physical time.  The physical observation grid differs from  the numerical integration grid: an ODE solver may use multiple internal function evaluations within each $I_k$ without changing the physical forecast horizon.

\paragraph{Action-conditioned forecasting.} The same construction extends directly to controlled dynamics by conditioning the velocity field on the action active at physical time. To make this precise, let $\va_{m+k}$ denote the action applied
during the transition from $\vx_{m+k}$ to
$\vx_{m+k+1}$.  Define the zero-order-hold control
$\va(s)
:=
\va_{m+k},$
$s\in I_k.$
The controlled VSF dynamics are

\begin{equation}
\frac{d\widehat{\vx}_s}{ds}
=
\vv_{{\theta}}
\bigl(
\widehat{\vx}_s,
s,
\vh,
\va(s),
\vz
\bigr).
\label{eq:vsf_action_conditioned_dynamics}
\end{equation}

During training, the aligned action is supplied to both the velocity field
and the local variational encoder.  In particular, the posterior becomes $q_{\boldsymbol{\phi}}
\bigl(
\vz
\mid
\vu_S,
\vc_S,
\va(S)
\bigr).$
The autonomous model is recovered by omitting the action argument.



For model-based planning, let $\va_{m:T-1}^{(i)}$ denote a candidate
action sequence.  For each candidate, draw $N_{\mathrm{ens}}$ independent
latents from $p(\vz)$ and generate terminal states
$\widehat{\vx}_{T}^{(i,j)}$.  The expected terminal cost is estimated by
\begin{equation}
\widehat{J}_{\vtheta}^{(i)}
:=
\frac{1}{N_{\mathrm{ens}}}
\sum_{j=1}^{N_{\mathrm{ens}}}
d_g
\bigl(
\widehat{\vx}_T^{(i,j)},
\vx_g
\bigr).
\label{eq:vsf_world_model_cost}
\end{equation}
Candidate sequences can then be optimized using CEM~\citep{rubinstein1999cross} or other planners.


